

\section{总体分析}

本文围绕零售销售额预测与数据分析这一核心任务，按照从数据探索到预测建模、再到业务应用的研究路径，递进式地解决四个问题。

首先，针对问题一进行数据基本统计与可视化分析，通过对9800条订单明细数据的全面描述性统计和时间序列可视化，建立对零售销售数据整体特征、时间趋势和维度分布的系统认知，为后续深入分析奠定数据基础。接着，针对问题二进行销售影响因素的相关性分析，采用单因素方差分析（ANOVA）和Spearman秩相关分析等方法，从产品、客户、区域、物流等多维度量化各因素对销售额的影响程度和统计显著性，识别出产品子类别和产品大类是驱动销售额差异的核心因素。然后，针对问题三构建销售额预测模型，基于2022-2024年训练数据，对比Holt-Winters指数平滑、随机森林、XGBoost等多种建模方法的预测性能，通过时间序列交叉验证和测试集评估筛选出最优模型，实现了对2025年月度销售额的有效预测。最后，针对问题四进行2026年未来销售额预测和经营策略建议，基于最优预测模型的滚动预测结果，结合影响因素分析的结论，从产品布局、区域营销、客户运营、物流履约和库存管理五个维度提出数据驱动的经营建议。

四个问题之间呈现层层递进的逻辑关系：问题一的数据探索为问题二的归因分析提供统计基础，问题二识别出的关键影响因素为问题三的特征工程提供变量选择依据，问题三的最优预测模型直接支撑问题四的未来预测，而问题四的策略建议又综合了问题二的归因结论和问题三的预测结果。整体研究框架从数据出发，经由分析和建模，最终回归到业务决策，形成了一个完整的数据驱动决策闭环。
